% ============================================================
% CAS/Python ενότητες — Κεφάλαιο 3: Γραμμικά Συστήματα
%
% QR code: qr_d7aaf59e2500.png  (→ latex/qrcodes/)
% URL: https://nikmatz.github.io/ELADIC_GR/code/ch03/
% ============================================================

\casactivbox%
  {Maxima: Γραμμικά Συστήματα — Gauss, Cramer, Ομογενή}%
  {%
    \label{cas:ch3_linear_systems}
    Εξετάστε τα παρακάτω τρία συστήματα $3\times 3$:
    \[
      (A)\;
      \begin{cases}
        2x + y - z = 8\\
       -3x - y + 2z = -11\\
        x - 2y + 2z = -3
      \end{cases}
      \qquad
      (B)\;
      \begin{cases}
        x + y = 3\\
        2x + 2y = 5
      \end{cases}
      \qquad
      (C)\;
      \begin{cases}
        x + 2y - z = 3\\
        2x + 4y - 2z = 6
      \end{cases}
    \]
    \begin{enumerate}[label=(\alph*),itemsep=2pt]
      \item Λύστε το σύστημα $(A)$ με \texttt{linsolve} και επαληθεύστε.
        Κατασκευάστε τον επαυξημένο πίνακα $[A|\mathbf{b}]$ και υπολογίστε
        \texttt{rref}.
      \item Για τα $(B)$ και $(C)$: υπολογίστε $\operatorname{rank}(A)$ και
        $\operatorname{rank}([A|\mathbf{b}])$. Εφαρμόστε το θεώρημα
        Rouché–Capelli και χαρακτηρίστε κάθε σύστημα.
      \item Λύστε το $2\times 2$ σύστημα $2x+y=4$, $5x+3y=7$ με τον
        κανόνα Cramer (αντικατάσταση στήλης, \texttt{determinant}).
      \item Βρείτε τον μηδενόχωρο (\texttt{nullspace}) του ομογενούς
        συστήματος $A\mathbf{x}=\mathbf{0}$ για
        $A=\bigl[\begin{smallmatrix}1&-2&1\\2&-3&1\\0&1&-1\end{smallmatrix}\bigr]$.
    \end{enumerate}
    \smallskip
    \textbf{Εντολές Maxima:}
    \texttt{linsolve([eq,...],[x,y,z])}, \texttt{rref(A)},
    \texttt{rank(A)}, \texttt{determinant(A)}, \texttt{nullspace(A)}
  }%
  {https://nikmatz.github.io/ELADIC_GR/code/ch03/}%
  {qr_d7aaf59e2500.png}

\subsection*{Δραστηριότητα Python}

\begin{pybox}{Python: Γραμμικά Συστήματα — NumPy, SymPy, Γεωμετρία}%
             {https://nikmatz.github.io/ELADIC_GR/code/ch03/}%
             {qr_d7aaf59e2500.png}
\label{py:ch3_linear_systems}
Χρησιμοποιήστε τα συστήματα $(A)$, $(B)$, $(C)$ της δραστηριότητας Maxima.
\begin{enumerate}[label=(\alph*),itemsep=2pt]
  \item Λύστε το $(A)$ με \texttt{np.linalg.solve}. Επαληθεύστε και
    εκτυπώστε την RREF με \texttt{sympy.Matrix.rref()}.
  \item Γράψτε συνάρτηση \texttt{classify\_system(A, b)} που ελέγχει
    τις τάξεις και επιστρέφει «μοναδική / άπειρες / ασύμβατο».
    Δοκιμάστε τη στα $(A)$, $(B)$, $(C)$.
  \item Για το $(C)$: εκτυπώστε την παραμετρική λύση με
    \texttt{sympy.linsolve}.
  \item Οπτικοποιήστε γεωμετρικά τους τρεις τύπους (τομή / παράλληλες /
    ταυτόσημες ευθείες) σε $2\times 2$ συστήματα με Matplotlib.
\end{enumerate}
\medskip
\textbf{Βιβλιοθήκες / Εντολές:}
\texttt{np.linalg.solve()}, \texttt{np.linalg.matrix\_rank()},
\texttt{sympy.Matrix.rref()}, \texttt{sympy.linsolve()},
\texttt{matplotlib.pyplot}
\end{pybox}
